%% toc: P15 157 157.5,157.6,157.7
% Proofs for the three-part decomposition of a real law, its regimes, two
% mixed laws, and the null modifications of a conditional-expectation version.
\section{Mixed laws and null modifications: proofs}
\label{sec:annex-p15-w02b-p15-002-mixed-laws}

Use the definitions of
Chapter~\ref{sec:p15-w02b-p15-002-mixed-laws}. Throughout, a Borel
measure means a measure on \((\mathbb R,\mathcal B(\mathbb R))\). Three
preliminary facts are used several times.

\emph{Fact~P1: sums, multiples and restrictions of measures.} Let
\(\nu_1,\nu_2\) be measures on a measurable space
\((S,\mathcal S)\), let \(c\geq0\) be finite and let \(B\in\mathcal S\).
Then \(\nu_1+\nu_2\), \(c\nu_1\) and \(\nu_1(\,\cdot\,\cap B)\) are
measures, and for every measurable \(g:S\to[0,\infty]\)
\[
 \begin{aligned}
 \int g\,d(\nu_1+\nu_2)&=\int g\,d\nu_1+\int g\,d\nu_2,\\
 \int g\,d(c\nu_1)&=c\int g\,d\nu_1,\\
 \int g\,d\nu_1(\,\cdot\,\cap B)&=\int_B g\,d\nu_1 .
 \end{aligned}
\]
\emph{Proof.} Each of the three set functions vanishes at
\(\varnothing\). For pairwise disjoint \(E_n\), countable additivity of
\(\nu_1(\,\cdot\,\cap B)\) is that of \(\nu_1\) applied to the disjoint
sets \(E_n\cap B\); that of \(c\nu_1\) follows because a nonnegative sum
is the increasing limit of its partial sums and multiplication by a
finite \(c\geq0\) commutes with finite sums and with increasing limits
(for \(c=0\) both sides are zero). For \(\nu_1+\nu_2\), the partial sums
of \(\sum_n(\nu_1(E_n)+\nu_2(E_n))\) are the sums of the partial sums of
the two series; both increase, so their limits add. For a nonnegative
simple \(s=\sum_ja_j\mathbf1_{A_j}\), the defining formula of
\cref{def:p15-simple-integral} gives the three identities directly; for
the restriction, \(\sum_ja_j\nu_1(A_j\cap B)\) is the simple integral of
\(s\mathbf1_B=\sum_ja_j\mathbf1_{A_j\cap B}\) by
\cref{thm:p15-simple-integral-laws}. For general \(g\), choose
nonnegative simple \(s_n\uparrow g\) (\cref{thm:p15-measurable-limits});
then \(s_n\mathbf1_B\uparrow g\mathbf1_B\), and monotone convergence
(\cref{thm:p15-monotone-integral}) applied to each measure passes the
three identities to the limit, since sums and finite nonnegative
multiples of increasing sequences in \([0,\infty]\) converge to the sums
and multiples of the limits.

\emph{Fact~P2: the distribution function of a finite Borel measure.}
Let \(\nu\) be a finite Borel measure and \(G(x)=\nu((-\infty,x])\).
Then \(0\leq G\leq\nu(\mathbb R)\), \(G\) is nondecreasing and right
continuous, \(G(x)\to0\) as \(x\to-\infty\),
\(G(x-)=\nu((-\infty,x))\), and \(G(x)-G(x-)=\nu(\{x\})\). In
particular \(G\) is continuous exactly when \(\nu\) is atomless.
\emph{Proof.} Monotonicity of \(\nu\) (\cref{thm:p15-measure-laws},
item~1) gives the bounds and that \(G\) is nondecreasing. For a
nondecreasing function, a one-sided limit at a point or at \(-\infty\)
equals the limit along any monotone sequence approaching it. The sets
\((-\infty,x+1/n]\) decrease to \((-\infty,x]\) and the sets
\((-\infty,-n]\) decrease to \(\varnothing\); continuity from above for
the finite measure \(\nu\) (\cref{thm:p15-measure-laws}, item~4) gives
right continuity and the limit zero. The sets \((-\infty,x-1/n]\)
increase to \((-\infty,x)\), so continuity from below gives
\(G(x-)=\nu((-\infty,x))\), and finite additivity gives
\(G(x)-G(x-)=\nu(\{x\})\). A nondecreasing right-continuous function is
continuous at \(x\) exactly when \(G(x-)=G(x)\).

\emph{Fact~P3: a squeeze integral.} Let \((\Omega,\mathcal F,m)\) be a
measure space, let \(V\) be a real measurable function, and let
\(0<c<\infty\) satisfy \(m(V\in(s,t])=t-s\) whenever \(0\leq s<t\leq c\).
Then
\[
 \int V\mathbf1_{\{0<V\leq c\}}\,dm=\frac{c^2}{2}.
\]
\emph{Proof.} Fix an integer \(n\geq1\) and put
\(J_k=((k-1)c/n,kc/n]\) for \(k=1,\ldots,n\). These intervals are
disjoint with union \((0,c]\). The simple functions
\[
 s_n=\sum_{k=1}^n\frac{(k-1)c}{n}\mathbf1_{\{V\in J_k\}},\qquad
 t_n=\sum_{k=1}^n\frac{kc}{n}\mathbf1_{\{V\in J_k\}}
\]
satisfy \(s_n\leq V\mathbf1_{\{0<V\leq c\}}\leq t_n\) at every point,
since on \(\{V\in J_k\}\) the value of \(V\) lies in \(J_k\) and off
\(\{0<V\leq c\}\) all three functions vanish. Each \(\{V\in J_k\}\) has
measure \(c/n\). Using \(\sum_{k=1}^nk=n(n+1)/2\) (true at \(n=1\), and
adding \(n+1\) to \(n(n+1)/2\) gives \((n+1)(n+2)/2\)), their integrals
are
\[
 \begin{aligned}
 \int s_n\,dm&=\frac{c^2}{n^2}\cdot\frac{(n-1)n}{2}
            =\frac{c^2}{2}\Bigl(1-\frac1n\Bigr),\\
 \int t_n\,dm&=\frac{c^2}{n^2}\cdot\frac{n(n+1)}{2}
            =\frac{c^2}{2}\Bigl(1+\frac1n\Bigr).
 \end{aligned}
\]
By monotonicity of the integral (\cref{thm:p15-monotone-integral}) the
integral in question lies between these numbers for every \(n\), so it
equals \(c^2/2\).

Two instances are used. For Lebesgue measure and \(V(x)=x\),
\(\lambda((s,t])=t-s\) by \cref{thm:p15-lebesgue-unit-interval}; since
a singleton is \(\lambda\)-null, part~(2) of
\cref{thm:p15-nonnegative-integral-laws} gives
\[
 \int_{(0,c)}x\,d\lambda(x)=\int_{(0,c]}x\,d\lambda(x)=\frac{c^2}{2}.
\]
On the uniform probability space, for \(0\leq s<t\leq1\),
\(\mathbb U(U\in(s,t])=\lambda((s,t]\cap(0,1))=t-s\), so for
\(0<c\leq1\)
\[
 \int_{\{U\leq c\}}U\,d\mathbb U=\frac{c^2}{2},\qquad
 \mathbb E[U]=\frac12,
\]
the second by taking \(c=1\), since \(0<U<1\) at every point.

\begin{annexitem}{annex:p15-w02b-p15-002-atomic-split}{thm:p15-w02b-p15-002-atomic-split}{Counting atoms and splitting them off}
\begin{proof}
(1) Fix \(k\geq1\) and let \(B_k=\{x:\nu(\{x\})\geq1/k\}\). If
\(x_1,\ldots,x_m\) are distinct points of \(B_k\), then finite
additivity and monotonicity give
\[
 \frac mk\leq\sum_{j=1}^m\nu(\{x_j\})=\nu(\{x_1,\ldots,x_m\})
        \leq\nu(\mathbb R),
\]
so \(m\leq k\,\nu(\mathbb R)\). Hence \(B_k\) is finite with at most
\(k\,\nu(\mathbb R)\) points: otherwise it would contain
\(m=\lfloor k\,\nu(\mathbb R)\rfloor+1\) distinct points. If
\(\nu(\{x\})>0\), then \(\nu(\{x\})\geq1/k\) for some integer \(k\).
Thus \(A_\nu=\bigcup_{k\geq1}B_k\) is a countable union of finite sets
and is countable. Being countable, it is Borel.

(2) By Fact~P1, \(\nu_{\mathrm d}=\nu(\,\cdot\,\cap A_\nu)\) and
\(\nu_{\mathrm c}=\nu(\,\cdot\,\cap(\mathbb R\setminus A_\nu))\) are
Borel measures. For Borel \(E\), the set \(E\cap A_\nu\) is countable
and is the disjoint union of its singletons, so countable additivity
(finite additivity when it is finite) and the fact that a nonnegative
sum does not depend on the enumeration of its index set
(\cref{def:p15-measure}) give
\(\nu(E\cap A_\nu)=\sum_{a\in E\cap A_\nu}\nu(\{a\})\). Since
\(\nu_{\mathrm d}(\mathbb R\setminus A_\nu)=\nu(\varnothing)=0\) and
\(A_\nu\) is countable, \(\nu_{\mathrm d}\) is purely atomic. For a
real \(x\), if \(x\in A_\nu\) then
\(\nu_{\mathrm c}(\{x\})=\nu(\varnothing)=0\), and if \(x\notin A_\nu\)
then \(\nu_{\mathrm c}(\{x\})=\nu(\{x\})=0\) by the definition of
\(A_\nu\). So \(\nu_{\mathrm c}\) is atomless. Finally
\(\nu(E)=\nu(E\cap A_\nu)+\nu(E\setminus A_\nu)\) by finite additivity,
which is \(\nu=\nu_{\mathrm d}+\nu_{\mathrm c}\).

(3) Let \(\nu=\nu_1+\nu_2\) with \(\nu_1(\mathbb R\setminus S)=0\) for a
countable \(S\) and \(\nu_2\) atomless. Both are finite, being at most
\(\nu\). For each \(x\),
\(\nu(\{x\})=\nu_1(\{x\})+\nu_2(\{x\})=\nu_1(\{x\})\), so \(\nu_1\) and
\(\nu\) have the same atoms. Every point of \(S\setminus A_\nu\) is
therefore \(\nu_1\)-null, and the countable set \(S\setminus A_\nu\) is
\(\nu_1\)-null by countable additivity. Hence
\[
 \nu_1(\mathbb R\setminus A_\nu)
   \leq\nu_1(\mathbb R\setminus S)+\nu_1(S\setminus A_\nu)=0
\]
by countable subadditivity (\cref{thm:p15-measure-laws}, item~2). Also
\(\nu_2(A_\nu)=\sum_{a\in A_\nu}\nu_2(\{a\})=0\). For Borel \(E\) it
follows that
\[
 \nu_{\mathrm d}(E)=\nu(E\cap A_\nu)
   =\nu_1(E\cap A_\nu)+\nu_2(E\cap A_\nu)=\nu_1(E\cap A_\nu)=\nu_1(E),
\]
the last step because
\(\nu_1(E\setminus A_\nu)\leq\nu_1(\mathbb R\setminus A_\nu)=0\).
Likewise
\[
 \nu_{\mathrm c}(E)=\nu(E\setminus A_\nu)
   =\nu_1(E\setminus A_\nu)+\nu_2(E\setminus A_\nu)
   =\nu_2(E\setminus A_\nu)=\nu_2(E),
\]
because \(\nu_2(E\cap A_\nu)\leq\nu_2(A_\nu)=0\).

(4) If \(\nu(\mathbb R\setminus A_\nu)=0\), then \(\nu\) is purely
atomic with the countable set \(S=A_\nu\) of part~(1). Conversely, if
\(\nu\) is purely atomic, apply part~(3) with \(\nu_1=\nu\) and
\(\nu_2=0\), which is atomless: \(\nu_{\mathrm c}=0\), and in particular
\(\nu(\mathbb R\setminus A_\nu)=\nu_{\mathrm c}(\mathbb R)=0\). Finally
\(\nu_{\mathrm c}=0\) holds exactly when
\(\nu_{\mathrm c}(\mathbb R)=\nu(\mathbb R\setminus A_\nu)=0\), because
\(0\leq\nu_{\mathrm c}(E)\leq\nu_{\mathrm c}(\mathbb R)\) for every
Borel \(E\).
\end{proof}
\end{annexitem}

\begin{annexitem}{annex:p15-w02b-p15-002-three-part}{thm:p15-w02b-p15-002-three-part}{The three parts: existence, uniqueness, weights and integrals}
\begin{proof}
On the probability space \((\mathbb R,\mathcal B(\mathbb R),\mu)\) the
identity map is a real random variable whose law is \(\mu\) and whose
distribution function is \(F\). By \cref{thm:p15-cdf-laws},
\(\mu(\{a\})=F(a)-F(a-)\) for every real \(a\), and the discontinuity
points of \(F\) are exactly the points of positive mass. Hence
\(A=A_\mu\), which is countable by part~(1) of
\cref{thm:p15-w02b-p15-002-atomic-split}.

(1) \emph{Existence.} Let \(\mu_{\mathrm d}\) and \(\mu_{\mathrm c}\)
be the measures of part~(2) of the lemma, so that
\(\mu=\mu_{\mathrm d}+\mu_{\mathrm c}\), \(\mu_{\mathrm d}\) is purely
atomic and \(\mu_{\mathrm c}\) is atomless, and
\[
 \mu_{\mathrm d}(E)=\sum_{a\in E\cap A}\mu(\{a\})
                   =\sum_{a\in E\cap A}\bigl(F(a)-F(a-)\bigr).
\]
The measure \(\mu_{\mathrm c}\) is finite, hence sigma-finite, and
\(\lambda\) is sigma-finite (\cref{thm:p15-lebesgue-unit-interval}). By
\cref{thm:p15-lebesgue-decomposition} there are measures
\(\mu_{\mathrm{ac}},\mu_{\mathrm{sc}}\) with
\(\mu_{\mathrm c}=\mu_{\mathrm{ac}}+\mu_{\mathrm{sc}}\),
\(\mu_{\mathrm{ac}}\ll\lambda\), \(\mu_{\mathrm{sc}}\perp\lambda\), and a
Borel \(g:\mathbb R\to[0,\infty]\), finite \(\lambda\)-almost
everywhere, with \(\mu_{\mathrm{ac}}(E)=\int_Eg\,d\lambda\). The set
\(\{g=\infty\}\) is Borel and \(\lambda\)-null; let \(f=g\) off it and
\(f=0\) on it. Then \(f\) is real, Borel and nonnegative, and
\(f\mathbf1_E=g\mathbf1_E\) \(\lambda\)-almost everywhere, so
\(\int_Ef\,d\lambda=\int_Eg\,d\lambda=\mu_{\mathrm{ac}}(E)\) by part~(2)
of \cref{thm:p15-nonnegative-integral-laws}. In particular
\(\int f\,d\lambda=\mu_{\mathrm{ac}}(\mathbb R)
\leq\mu_{\mathrm c}(\mathbb R)\leq\mu(\mathbb R)=1\). For every real
\(x\), \(\mu_{\mathrm{sc}}(\{x\})\leq\mu_{\mathrm c}(\{x\})=0\), so
\(\mu_{\mathrm{sc}}\) is atomless and singular, that is, continuous
singular. Thus \(\mu=\mu_{\mathrm d}+\mu_{\mathrm{ac}}+\mu_{\mathrm{sc}}\)
with the three required types.

(2) \emph{Uniqueness.} Let \(\mu=\nu_1+\nu_2+\nu_3\) as in the
statement. Each \(\nu_i\) is at most \(\mu\), hence finite. Since
\(\lambda(\{x\})=0\) and \(\nu_2\ll\lambda\), the measure \(\nu_2\) is
atomless; \(\nu_3\) is atomless by definition; so \(\nu_2+\nu_3\) is an
atomless measure (Fact~P1). Part~(3) of
\cref{thm:p15-w02b-p15-002-atomic-split}, applied to
\(\mu=\nu_1+(\nu_2+\nu_3)\), gives \(\nu_1=\mu_{\mathrm d}\) and
\(\nu_2+\nu_3=\mu_{\mathrm c}\). Now
\(\mu_{\mathrm c}=\nu_2+\nu_3=\mu_{\mathrm{ac}}+\mu_{\mathrm{sc}}\) are
two decompositions of the finite measure \(\mu_{\mathrm c}\) into a part
absolutely continuous with respect to \(\lambda\) and a part singular
with respect to \(\lambda\). The uniqueness clause of
\cref{thm:p15-lebesgue-decomposition}, with the sigma-finite measures
\(\mu_{\mathrm c}\) and \(\lambda\), gives \(\nu_2=\mu_{\mathrm{ac}}\)
and \(\nu_3=\mu_{\mathrm{sc}}\). If \(f\) and \(\tilde f\) are Borel
densities of \(\mu_{\mathrm{ac}}\) with respect to \(\lambda\), they
agree \(\lambda\)-almost everywhere by \cref{thm:p15-radon-nikodym},
applied to \(\mu_{\mathrm{ac}}\ll\lambda\).

(3) \emph{Distribution functions.} For every real \(x\),
\[
 F_{\mathrm d}(x)=\mu_{\mathrm d}((-\infty,x]),\quad
 F_{\mathrm{ac}}(x)=\mu_{\mathrm{ac}}((-\infty,x]),\quad
 F_{\mathrm{sc}}(x)=\mu_{\mathrm{sc}}((-\infty,x]):
\]
the first by the formula in part~(1), the second by the density
representation, the third by definition. Adding the three and using
part~(1) gives
\[
 F_{\mathrm d}(x)+F_{\mathrm{ac}}(x)+F_{\mathrm{sc}}(x)
   =\mu((-\infty,x])=F(x).
\]
Fact~P2, applied to the three finite measures, shows that each of the
three functions is bounded, nondecreasing and right continuous with
limit zero at \(-\infty\), and that its jump at \(x\) is the mass of the
corresponding measure at \(\{x\}\). The measures \(\mu_{\mathrm{ac}}\)
(as shown in part~(2) for any part absolutely continuous with respect to
\(\lambda\)) and \(\mu_{\mathrm{sc}}\) are atomless, so
\(F_{\mathrm{ac}}\) and \(F_{\mathrm{sc}}\) are continuous. For
\(F_{\mathrm d}\), the jump at \(x\) is
\(\mu_{\mathrm d}(\{x\})=\mu(\{x\}\cap A)\); if \(x\in A\) this is
\(\mu(\{x\})\), and if \(x\notin A\) it is \(0=\mu(\{x\})\) because
\(A=A_\mu\). In both cases it equals \(F(x)-F(x-)\).

For the uniqueness clause, let \(F=G_1+G_2+G_3\) as in the statement,
and let \(\mu_{G_i}\) be the Stieltjes measures of
\cref{thm:p15-stieltjes-construction}. Since \(G_i\) is nondecreasing
with limit zero at \(-\infty\), its infimum \(G_i(-\infty)\) is zero, so
\(\mu_{G_i}((-\infty,x])=G_i(x)-G_i(-\infty)=G_i(x)\), and
\(\mu_{G_i}(\mathbb R)=\sup G_i<\infty\). By Fact~P1,
\(\nu=\mu_{G_1}+\mu_{G_2}+\mu_{G_3}\) is a Borel measure with
\(\nu((-\infty,x])=F(x)\) for every real \(x\). The sets
\((-\infty,n]\) increase to \(\mathbb R\), so continuity from below and
\(F(n)\to1\) (\cref{thm:p15-cdf-laws}) give \(\nu(\mathbb R)=1\). The
identity maps on \((\mathbb R,\mathcal B(\mathbb R),\nu)\) and on
\((\mathbb R,\mathcal B(\mathbb R),\mu)\) are real random variables with
laws \(\nu\) and \(\mu\) and the same distribution function \(F\); by
\cref{thm:p15-cdf-laws} they have the same law, so \(\nu=\mu\). The
hypotheses on the types of \(\mu_{G_1},\mu_{G_2},\mu_{G_3}\) and part~(2)
give \(\mu_{G_1}=\mu_{\mathrm d}\), \(\mu_{G_2}=\mu_{\mathrm{ac}}\) and
\(\mu_{G_3}=\mu_{\mathrm{sc}}\). Evaluating on \((-\infty,x]\) gives
\(G_1=F_{\mathrm d}\), \(G_2=F_{\mathrm{ac}}\) and
\(G_3=F_{\mathrm{sc}}\).

(4) \emph{Mixture form.} The weights are values of measures, hence
nonnegative, and by part~(1)
\(p_{\mathrm d}+p_{\mathrm{ac}}+p_{\mathrm{sc}}=\mu(\mathbb R)=1\). Let
\(\mu_\bullet\) be one of the three parts with \(p_\bullet>0\). By
Fact~P1, \(\mu_\bullet/p_\bullet\) is a measure; its total mass is one.
Multiplying a measure by a positive constant leaves its null sets, the
sets on which it is concentrated and its atomless property unchanged.
Hence \(\mu_\bullet/p_\bullet\) is purely atomic, absolutely continuous
or continuous singular exactly as \(\mu_\bullet\) is, and its
distribution function is \(\mu_\bullet((-\infty,x])/p_\bullet
=F_\bullet(x)/p_\bullet\). If \(p_\bullet=0\), then
\(0\leq\mu_\bullet(E)\leq\mu_\bullet(\mathbb R)=0\) for all Borel \(E\),
so \(\mu_\bullet\) is the zero measure and \(F_\bullet=0\). Dropping
these zero terms from \(F=F_{\mathrm d}+F_{\mathrm{ac}}+F_{\mathrm{sc}}\)
and writing each remaining \(F_\bullet\) as
\(p_\bullet\,(F_\bullet/p_\bullet)\) gives the mixture.

(5) \emph{Integration.} Let \(g:\mathbb R\to[0,\infty]\) be Borel. By
part~(1) and Fact~P1,
\[
 \int g\,d\mu=\int g\,d\mu_{\mathrm d}+\int g\,d\mu_{\mathrm{ac}}
              +\int g\,d\mu_{\mathrm{sc}} .
\]
For the first term, Fact~P1 gives
\(\int g\,d\mu_{\mathrm d}=\int g\mathbf1_A\,d\mu\). If \(A\) is empty,
this is zero, the empty sum. Otherwise list \(A\) without repetition as
\(a_1,a_2,\ldots\) (a finite list when \(A\) is finite). At every point,
\(g\mathbf1_A=\sum_ng(a_n)\mathbf1_{\{a_n\}}\), since at most one term
is nonzero. For a constant \(c\in[0,\infty)\),
\(\int c\mathbf1_{\{a\}}\,d\mu=c\,\mu(\{a\})\)
(\cref{thm:p15-simple-integral-laws}); for \(c=\infty\), monotone
convergence applied to \(k\mathbf1_{\{a\}}\uparrow\infty\cdot\mathbf1_{\{a\}}\)
gives \(\lim_kk\,\mu(\{a\})=\infty=\infty\cdot\mu(\{a\})\), as
\(\mu(\{a\})>0\) for \(a\in A\). Termwise integration of a nonnegative
series (\cref{thm:p15-nonnegative-integral-laws}) and independence of a
nonnegative sum from the enumeration then give
\[
 \int g\,d\mu_{\mathrm d}=\sum_ng(a_n)\mu(\{a_n\})
     =\sum_{a\in A}g(a)\bigl(F(a)-F(a-)\bigr).
\]
For the second term, \(\mu_{\mathrm{ac}}\ll\lambda\), both measures are
sigma-finite and \(f\) is a density of \(\mu_{\mathrm{ac}}\), so item~1
of \cref{thm:p15-density-calculus} gives
\(\int g\,d\mu_{\mathrm{ac}}=\int gf\,d\lambda\). This proves the
nonnegative identity.

Now let \(g\) be real Borel with \(\int|g|\,d\mu<\infty\). Applying the
nonnegative identity to \(|g|\) shows that
\(\sum_{a\in A}|g(a)|(F(a)-F(a-))\), \(\int|g|f\,d\lambda\) and
\(\int|g|\,d\mu_{\mathrm{sc}}\) are finite, since their sum is
\(\int|g|\,d\mu\). Applying it to \(g^+\) and to \(g^-\) gives two
identities between finite numbers. In any enumeration of \(A\), the
partial sums of \(\sum_ng(a_n)\mu(\{a_n\})\) are differences of the
partial sums of the convergent nonnegative series for \(g^+\) and
\(g^-\); hence the series converges, absolutely by the bound for \(|g|\),
to the difference of their sums, whatever the enumeration. Since \(f\) is
real and nonnegative, \((gf)^+=g^+f\) and \((gf)^-=g^-f\), so
\(\int gf\,d\lambda=\int g^+f\,d\lambda-\int g^-f\,d\lambda\) with both
terms finite; likewise
\(\int g\,d\mu_{\mathrm{sc}}=\int g^+\,d\mu_{\mathrm{sc}}
-\int g^-\,d\mu_{\mathrm{sc}}\) and
\(\int g\,d\mu=\int g^+\,d\mu-\int g^-\,d\mu\)
(\cref{def:p15-integrable-functions}). Subtracting the identity for
\(g^-\) from the identity for \(g^+\) and regrouping the finite terms
gives the identity for \(g\).
\end{proof}
\end{annexitem}

\begin{annexitem}{annex:p15-w02b-p15-002-regimes}{thm:p15-w02b-p15-002-regimes}{Reading each regime off the three parts}
\begin{proof}
Throughout, use the parts and weights of
\cref{thm:p15-w02b-p15-002-three-part}. The zero measure is purely
atomic (it is concentrated on the empty set), absolutely continuous, and
continuous singular (take the null set \(\varnothing\) in
\cref{def:p15-singularity}).

(1) A probability law is discrete when it gives probability one to a
countable Borel set, which is the purely atomic property. By part~(4) of
\cref{thm:p15-w02b-p15-002-atomic-split}, \(\mu\) is purely atomic
exactly when \(\mu_{\mathrm c}=0\), that is, when
\(\mu=\mu_{\mathrm d}\). Moreover \(\mu=\mu_{\mathrm d}\) holds exactly
when \(\mu(\mathbb R\setminus A)=\mu_{\mathrm c}(\mathbb R)=0\), which,
since \(\mu(\mathbb R)=1\), holds exactly when
\(\mu_{\mathrm d}(\mathbb R)=\mu(A)=1\); and
\(\mu_{\mathrm d}(\mathbb R)=\sum_{a\in A}(F(a)-F(a-))\) by part~(1) of
the theorem.

(2) By definition \(A\) is the set of discontinuity points of \(F\), so
\(F\) is continuous exactly when \(A=\varnothing\). If
\(A=\varnothing\), then \(\mu_{\mathrm d}(E)=\mu(E\cap\varnothing)=0\)
for every \(E\). If \(A\) contains a point \(a\), then
\(\mu_{\mathrm d}(\{a\})=\mu(\{a\})>0\), so \(\mu_{\mathrm d}\neq0\).

(3) Suppose that \(\mu(E)=\int_Eh\,d\lambda\) for every Borel \(E\), for
some Borel \(h\geq0\) (\cref{def:p15-probability-density}). If
\(\lambda(E)=0\), then \(h\mathbf1_E=0\) \(\lambda\)-almost everywhere,
so \(\mu(E)=0\) by part~(1) of \cref{thm:p15-nonnegative-integral-laws};
thus \(\mu\ll\lambda\). Then \(\mu=0+\mu+0\) is a decomposition with a
purely atomic, an absolutely continuous and a continuous singular term,
and part~(2) of the theorem gives \(\mu=\mu_{\mathrm{ac}}\). Conversely,
if \(\mu=\mu_{\mathrm{ac}}\), then \(\mu(E)=\int_Ef\,d\lambda\) with the
density \(f\) of part~(1) of the theorem.

(4) Suppose \(\mu\perp\lambda\): some Borel \(N\) has \(\lambda(N)=0\)
and \(\mu(\mathbb R\setminus N)=0\). The measure \(\mu_{\mathrm c}\) of
\cref{thm:p15-w02b-p15-002-atomic-split} is atomless and satisfies
\(\mu_{\mathrm c}(\mathbb R\setminus N)\leq\mu(\mathbb R\setminus N)=0\),
so it is continuous singular. Hence
\(\mu=\mu_{\mathrm d}+0+\mu_{\mathrm c}\) is a decomposition of the
three types, and part~(2) of the theorem gives \(\mu_{\mathrm{ac}}=0\).
Conversely, suppose \(\mu_{\mathrm{ac}}=0\), so
\(\mu=\mu_{\mathrm d}+\mu_{\mathrm{sc}}\). Choose a Borel \(N_1\) with
\(\lambda(N_1)=0\) and \(\mu_{\mathrm{sc}}(\mathbb R\setminus N_1)=0\),
and put \(N=N_1\cup A\). The countable set \(A\) is \(\lambda\)-null, so
\(\lambda(N)=0\), and
\[
 \mu(\mathbb R\setminus N)
   \leq\mu_{\mathrm d}(\mathbb R\setminus A)
       +\mu_{\mathrm{sc}}(\mathbb R\setminus N_1)=0 .
\]
Thus \(\mu\perp\lambda\).

(5) By \cref{thm:p15-cantor-singular-law}, \(\eta\) is a Borel
probability law, every singleton is \(\eta\)-null, \(\eta(C)=1\) and
\(\lambda(C)=0\). The set \(C\) is Borel: each \(C_n\) is a finite union
of closed intervals and \(C\) is their countable intersection. Since
\(\eta(\mathbb R\setminus C)=1-\eta(C)=0\), the set \(N=C\) witnesses
\(\eta\perp\lambda\). Hence \(\eta\) is continuous singular, and
\(\eta=0+0+\eta\) is a decomposition of the three types; part~(2) of the
theorem gives \(\eta_{\mathrm d}=\eta_{\mathrm{ac}}=0\) and
\(\eta_{\mathrm{sc}}=\eta\).

\emph{Mixed laws.} The three classes named are those of the laws with
\(\mu=\mu_{\mathrm d}\) (item~1), with \(\mu=\mu_{\mathrm{ac}}\)
(item~3), and with \(\mu=\mu_{\mathrm{sc}}\), the class to which item~5
assigns \(\eta\). If \(\mu=\mu_\bullet\) for one of the three parts, the
other two parts have total mass
\(\mu(\mathbb R)-\mu_\bullet(\mathbb R)=0\), so their weights vanish and
at most one weight is positive. A mixed law has at least two positive
weights, so it lies in none of the three classes.
\end{proof}
\end{annexitem}

\begin{annexitem}{annex:p15-w02b-p15-002-mixed-examples}{thm:p15-w02b-p15-002-mixed-examples}{The censored uniform law and the atomic, Lebesgue and Cantor mixture}
\begin{proof}
(1) \emph{Measurability and distribution function.} On
\(\Omega_U=(0,1)\) the function \(X=\min(U,1/2)\) is a real random
variable by \cref{thm:p15-measurable-operations}. Since \(U>0\), also
\(X>0\) at every point, so \(\{X\leq x\}=\varnothing\) and
\(F_X(x)=0\) for \(x<0\), and also \(F_X(0)=0\). For \(0<x<1/2\),
\(\min(U,1/2)\leq x\) holds exactly when \(U\leq x\), because
\(1/2>x\); so \(F_X(x)=\mathbb U(U\leq x)=x\) by
\cref{thm:p15-lebesgue-unit-interval}. For \(x\geq1/2\), the event is
all of \(\Omega_U\), so \(F_X(x)=1\). This is the displayed formula. The
function is constant on \((-\infty,0)\) and on \([1/2,\infty)\), equal
to \(x\) on \([0,1/2)\), and continuous at \(0\); its only discontinuity
is at \(1/2\), where \(F_X(1/2-)=\sup_{t<1/2}F_X(t)=1/2\) and
\(F_X(1/2)=1\), a jump of \(1/2\).

\emph{The parts.} Put \(\nu_1=\tfrac12\delta_{1/2}\),
\(\nu_2(E)=\lambda(E\cap(0,1/2))\) and \(\nu_3=0\). By Fact~P1,
\(\nu_1+\nu_2\) is a Borel measure, and
\[
 (\nu_1+\nu_2)((-\infty,x])=
 \tfrac12\mathbf1_{[1/2,\infty)}(x)+\lambda((-\infty,x]\cap(0,1/2)),
\]
which is \(0\) for \(x\leq0\), \(x\) for \(0<x<1/2\) and
\(\tfrac12+\tfrac12=1\) for \(x\geq1/2\). This is \(F_X\), so
\(\nu_1+\nu_2\) is a probability measure, and by
\cref{thm:p15-cdf-laws} (comparing the identity variable on
\((\mathbb R,\mathcal B(\mathbb R),\nu_1+\nu_2)\) with \(X\)) the law
\(\mu\) of \(X\) equals \(\nu_1+\nu_2+\nu_3\). Here \(\nu_1\) is
concentrated on \(\{1/2\}\), so it is purely atomic; \(\nu_2(E)
=\int\mathbf1_E\mathbf1_{(0,1/2)}\,d\lambda=\int_E\mathbf1_{(0,1/2)}\,d\lambda\)
by \cref{thm:p15-simple-integral-laws}, so \(\nu_2\) has density
\(\mathbf1_{(0,1/2)}\) and is absolutely continuous by
\cref{thm:p15-radon-nikodym}; and \(\nu_3=0\) is continuous singular.
Part~(2) of \cref{thm:p15-w02b-p15-002-three-part} gives
\(\mu_{\mathrm d}=\tfrac12\delta_{1/2}\),
\(\mu_{\mathrm{ac}}=\nu_2\) and \(\mu_{\mathrm{sc}}=0\).

\emph{Mixture form.} The weights are
\(p_{\mathrm d}=\tfrac12\), \(p_{\mathrm{ac}}=\lambda((0,1/2))=\tfrac12\)
and \(p_{\mathrm{sc}}=0\). By part~(4) of the theorem the normalized
parts are \(\delta_{1/2}\) and \(2\nu_2\), and
\(2\nu_2(E)=\int_E2\cdot\mathbf1_{(0,1/2)}\,d\lambda\) by Fact~P1, so
\(2\cdot\mathbf1_{(0,1/2)}\) is a density of the second; its total mass
is \(2\lambda((0,1/2))=1\).

\emph{The mean.} Let \(g(x)=\max(x,0)\), a nonnegative Borel function.
Since \(X>0\), \(X=g(X)\), so
\(\mathbb E[X]=\int g\,d\mu\) by \cref{thm:p15-expectation-law-integral}.
Part~(5) of the theorem, with \(A=\{1/2\}\), density
\(\mathbf1_{(0,1/2)}\) and \(\mu_{\mathrm{sc}}=0\) (all of whose
integrals are zero, every simple integral being zero), gives
\[
 \mathbb E[X]=g(\tfrac12)\cdot\tfrac12
   +\int g\,\mathbf1_{(0,1/2)}\,d\lambda
 =\tfrac14+\int_{(0,1/2)}x\,d\lambda(x)=\tfrac14+\tfrac18=\tfrac38,
\]
using Fact~P3 with \(c=1/2\) and \(g(x)=x\) on \((0,1/2)\). Directly on
the uniform space, \(\min(U,1/2)=U\mathbf1_{\{U\leq1/2\}}
+\tfrac12\mathbf1_{\{U>1/2\}}\) at every point, so additivity of
nonnegative integrals, Fact~P3 with \(c=1/2\) and
\(\mathbb U(U>1/2)=\tfrac12\) give
\(\mathbb E[\min(U,1/2)]=\tfrac18+\tfrac12\cdot\tfrac12=\tfrac38\).

(2) \emph{The law \(\rho\).} By Fact~P1, \(\rho\) is a Borel measure,
and \(\rho(\mathbb R)=\tfrac14+\tfrac12\lambda((0,1))+\tfrac14
\eta(\mathbb R)=\tfrac14+\tfrac12+\tfrac14=1\). Its three terms have the
required types: \(\tfrac14\delta_2\) is concentrated on \(\{2\}\);
\(\tfrac12\lambda(\,\cdot\,\cap(0,1))\) has density
\(\tfrac12\mathbf1_{(0,1)}\) by \cref{thm:p15-simple-integral-laws} and
Fact~P1, hence is absolutely continuous; and \(\tfrac14\eta\) is
continuous singular by item~5 of \cref{thm:p15-w02b-p15-002-regimes},
whose proof shows that \(\eta\) is atomless and concentrated on the
\(\lambda\)-null set \(C\), properties kept by a positive multiple.
Part~(2) of \cref{thm:p15-w02b-p15-002-three-part} gives
\(\rho_{\mathrm d}=\tfrac14\delta_2\),
\(\rho_{\mathrm{ac}}=\tfrac12\lambda(\,\cdot\,\cap(0,1))\) and
\(\rho_{\mathrm{sc}}=\tfrac14\eta\).

\emph{Distribution function.} Evaluate the three terms on
\((-\infty,x]\). First, the point mass at \(2\) gives
\(\mathbf1_{[2,\infty)}(x)\). Next,
\((-\infty,x]\cap(0,1)\) is empty for \(x\leq0\), equals \((0,x]\) for
\(0<x<1\) and \((0,1)\) for \(x\geq1\), so its Lebesgue measure is
\(\min(\max(x,0),1)\). Finally \(\eta((-\infty,x])=F_\eta(x)\). This is the
displayed formula for \(F_\rho\). By Fact~P2 the jump of \(F_\rho\) at
\(x\) is
\[
 \rho(\{x\})=\tfrac14\delta_2(\{x\})+\tfrac12\lambda(\{x\}\cap(0,1))
             +\tfrac14\eta(\{x\})=\tfrac14\mathbf1_{\{x=2\}},
\]
because singletons are \(\lambda\)-null and \(\eta\)-null. So \(F_\rho\)
has exactly one jump, of size \(1/4\) at \(2\). For
\(1/3\leq x\leq2/3\), \(\mathbf1_{[2,\infty)}(x)=0\),
\(\min(\max(x,0),1)=x\) and \(F_\eta(x)=1/2\) by
\cref{thm:p15-cantor-singular-law}, so \(F_\rho(x)=x/2+1/8\). Since
\(F_\eta(0)=0\) and \(F_\eta(1)=1\) by the same theorem,
\[
 F_\rho(1)-F_\rho(0)=\bigl(0+\tfrac12+\tfrac14\bigr)-\bigl(0+0+0\bigr)
                    =\tfrac34,
\]
while the density part gives
\(\rho_{\mathrm{ac}}([0,1])=\tfrac12\lambda([0,1]\cap(0,1))=\tfrac12\),
and also \(\rho_{\mathrm{ac}}((0,1])=\tfrac12\).

\emph{The Cantor set.} \(\lambda(C)=0\) by
\cref{thm:p15-cantor-singular-law}. Since \(C\subseteq C_1\subseteq
[0,1]\), the point \(2\) is not in \(C\), and
\(\lambda(C\cap(0,1))\leq\lambda(C)=0\). Hence
\[
 \rho(C)=\tfrac14\delta_2(C)+\tfrac12\lambda(C\cap(0,1))
         +\tfrac14\eta(C)=0+0+\tfrac14\cdot1=\tfrac14 .
\]
\end{proof}
\end{annexitem}

\begin{annexitem}{annex:p15-w02b-p15-002-null-modification}{thm:p15-w02b-p15-002-null-modification}{Measurability decides null modifications}
\begin{proof}
By \cref{thm:p15-conditional-existence}, \(Z\) is integrable.

(1) If \(Z'\) is a version, it is \(\mathcal G\)-measurable by
\cref{def:p15-general-conditional-expectation}. Conversely, let \(Z'\)
be \(\mathcal G\)-measurable with \(Z'=Z\) almost surely, so the event
\(D=\{Z'\neq Z\}\in\mathcal F_0\) is null. For \(A\in\mathcal G\), the
functions \(Z'\mathbf1_A\) and \(Z\mathbf1_A\) agree outside \(D\), and
\(Z\mathbf1_A\) is integrable because \(|Z\mathbf1_A|\leq|Z|\). By
\cref{thm:p15-signed-integral-laws}, measurable almost-everywhere
replacement preserves integrability and the integral, so the event
integrals of \(Z'\) are defined and
\(\int_AZ'\,dP=\int_AZ\,dP=\int_AX\,dP\). Hence \(Z'\) is a version.
This is the case of part~(ii) of
\cref{thm:p15-w02b-p15-008-null-modification} in which \(Z'\) is
\(\mathcal F_0\)-measurable, with \(D'=D\).

(2) The sets \(N\) and \(N^c\) belong to \(\mathcal G\), so their
indicators are \(\mathcal G\)-measurable, and
\(Z''=Z\mathbf1_{N^c}+W\mathbf1_N\) is a real \(\mathcal G\)-measurable
function by \cref{thm:p15-measurable-operations}, applied on
\((\Omega,\mathcal G)\). It agrees with \(Z\) outside the null set \(N\),
so it equals \(Z\) almost surely, and it is \(\mathcal F_0\)-measurable
because \(\mathcal G\subseteq\mathcal F_0\). By item~1 it is a version.
This is part~(i) of \cref{thm:p15-w02b-p15-008-null-modification}.

(3) The function \(Z+\mathbf1_N\) is real, \(\mathcal F_0\)-measurable,
and differs from \(Z\) only on the null event \(N\). By item~1 it is a
version exactly when it is \(\mathcal G\)-measurable. If
\(N\in\mathcal G\), then \(\mathbf1_N\) and hence \(Z+\mathbf1_N\) are
\(\mathcal G\)-measurable. Conversely, if \(Z+\mathbf1_N\) is
\(\mathcal G\)-measurable, then so is the difference
\(\mathbf1_N=(Z+\mathbf1_N)-Z\) of two real \(\mathcal G\)-measurable
functions, and therefore \(N=\{\mathbf1_N=1\}\in\mathcal G\).

(4) Let \(Z'\) be real and \(\mathcal F_0\)-measurable with \(Z'=Z\)
almost surely, and let \(D=\{Z'\neq Z\}\), an event of probability zero.
By hypothesis \(D\in\mathcal G\). For a Borel set \(B\),
\[
 \{Z'\in B\}=\bigl(\{Z\in B\}\cap D^c\bigr)\cup
             \bigl(\{Z'\in B\}\cap D\bigr).
\]
The first set is in \(\mathcal G\) because \(Z\) is
\(\mathcal G\)-measurable and \(D^c\in\mathcal G\). The second belongs
to \(\mathcal F_0\) and has probability at most \(P(D)=0\), so it is a
\(P\)-null set of \(\mathcal F_0\) and belongs to \(\mathcal G\) by
hypothesis. Thus \(Z'\) is \(\mathcal G\)-measurable, and item~1 shows
that it is a version.

(5) Let \(\mathcal G=\{\varnothing,\Omega\}\). Since \(P(\Omega)=1\),
the set \(\Omega\) is nonempty. Let \(Z'\) be any version and fix
\(\omega_0\in\Omega\). The set \(\{Z'=Z'(\omega_0)\}\) is the inverse
image of a Borel singleton, so it belongs to \(\mathcal G\); it contains
\(\omega_0\), so it is \(\Omega\). Thus \(Z'\) is a constant \(c\), and
the defining identity with \(A=\Omega\) gives
\(c=\int_\Omega Z'\,dP=\int_\Omega X\,dP=\mathbb E[X]\). Conversely the
constant \(\mathbb E[X]\) is \(\mathcal G\)-measurable, and its event
integrals over \(\varnothing\) and \(\Omega\) are \(0\) and
\(\mathbb E[X]\), those of \(X\). So the constant \(\mathbb E[X]\) is a
version and every version equals it at every point. The only member of
\(\mathcal G\) with probability zero is \(\varnothing\), since
\(P(\Omega)=1\). If \(N\in\mathcal F_0\) is nonempty with \(P(N)=0\),
then \(N\notin\mathcal G\), so \(\mathbb E[X]+\mathbf1_N\), which equals
the version outside the null event \(N\), is not a version by item~3.

For the uniform probability space, \(X=U\) is integrable since
\(0<U<1\), and \(\mathbb E[U]=1/2\) by Fact~P3. The singleton
\(N=\{1/2\}\) belongs to \(\mathcal F_U\) and
\(\mathbb U(N)=\lambda(\{1/2\})=0\) by
\cref{thm:p15-lebesgue-unit-interval}. So the only version is the
constant \(1/2\), and the function equal to \(1/2\) except for the value
\(3/2\) at \(u=1/2\) is \(\tfrac12+\mathbf1_N\), which is not a version.

(6) The set \(L=(0,1/2]=(0,1/2]\cap(0,1)\) belongs to \(\mathcal F_U\),
and the four-member family \(\mathcal G\) contains \(\Omega_U\) and is
closed under complements (\(\varnothing\leftrightarrow\Omega_U\),
\(L\leftrightarrow L^c\)) and under arbitrary unions of its members,
each such union again being one of the four sets; so it is a
sub-sigma-algebra of \(\mathcal F_U\). The function
\(Z=\tfrac14\mathbf1_L+\tfrac34\mathbf1_{L^c}\) is
\(\mathcal G\)-measurable: the inverse image of any Borel set is
\(\varnothing\), \(L\), \(L^c\) or \(\Omega_U\), according to which of
\(1/4\) and \(3/4\) it contains. Its event integrals are
\(\int_LZ\,d\mathbb U=\tfrac14\mathbb U(L)=\tfrac18\) and
\(\int_{L^c}Z\,d\mathbb U=\tfrac34\mathbb U(L^c)=\tfrac38\). For \(U\),
Fact~P3 with \(c=1/2\) gives \(\int_LU\,d\mathbb U=\tfrac18\), and
additivity over the disjoint union \(\Omega_U=L\cup L^c\)
(\cref{thm:p15-nonnegative-integral-laws}, item~4) gives
\(\int_{L^c}U\,d\mathbb U=\mathbb E[U]-\tfrac18=\tfrac38\). The
integrals over \(\varnothing\) are zero and those over \(\Omega_U\) are
\(\tfrac18+\tfrac38=\tfrac12=\mathbb E[U]\). So \(Z\) is a version.

Let \(Z'\) be any version. For \(\omega_0\in L\), the set
\(\{Z'=Z'(\omega_0)\}\) is a member of \(\mathcal G\) containing
\(\omega_0\), hence it is \(L\) or \(\Omega_U\) and contains \(L\); so
\(Z'\) is constant on \(L\), with value \(a\) say. In the same way
\(Z'\) is constant on \(L^c\), with value \(b\). The defining identity on
\(L\) gives \(a/2=1/8\), and on \(L^c\) gives \(b/2=3/8\). Hence
\(a=1/4\), \(b=3/4\) and \(Z'=Z\) at every point. A function that
agrees with \(Z\) except at \(u=1/4\), where its value differs from
\(1/4\), equals \(Z\) outside the \(\mathbb U\)-null set \(\{1/4\}\),
but it is not \(Z\) at every point and so is not a version. Equivalently,
it is not constant on \(L\) and therefore not \(\mathcal G\)-measurable,
in accordance with item~3: \(\{1/4\}\notin\mathcal G\).
\end{proof}
\end{annexitem}
